\begin{theorem}
Let $1<p<\infty$ with $p\neq2$, let $k\ge1$ be an integer and
$0<\alpha\le1$ with $k+\alpha$ equal to the right-hand side of (9), as in
\noderef{KAlphaExp}, and put $q=\frac{2}{2-\alpha}$.  Let
$U\subseteq\mathbb{C}$ be open with $0\in U$, let $H=H_p$ be the extremal
function (54) of \noderef{ExtremalH}, let $d=d(p)$ be the exponent of
\noderef{DExp}, and let $f\colon\mathbb{C}\to\mathbb{C}$ satisfy
\begin{enumerate}
\item[(H1)] $f$ is $C^\infty$ on $\mathbb{C}-\{0\}$;
\item[(H2)] $f(H(\xi))=\xi$ for every $\xi\in\mathbb{C}$, and $H(f(z))=z$ for
  every $z\in\mathbb{C}$;
\item[(H3)] $f(tz)=t^{d}f(z)$ for every $t\ge0$ and every $z\in\mathbb{C}$;
\item[(H4)] $D^\nu f(tz)=t^{\,d-|\nu|}D^\nu f(z)$ for every $t>0$, every
  $z\neq0$ and every iterated partial derivative $D^\nu$ of
  \noderef{PartialWord}.
\end{enumerate}
Then
\[
  f\notin C^{k-1+\alpha}_{\mathrm{loc}}(U)
  \quad\text{and}\quad
  f\notin W^{k+1,q}_{\mathrm{loc}}(U),
\]
in the senses of \noderef{MemHolderSmallOn} and \noderef{MemWLocOn}; here
$k\ge1$, so $k-1$ is the ordinary predecessor.

This is (56) of the paper, the heart of the optimality half of Theorem 1.  Two
features of the statement are forced.  The two non-memberships are asserted on
an arbitrary open $U$ containing the singular point, not on $\mathbb{C}$: a
failure of the little-o condition, or of the integrability, is witnessed by a
compact subset of the domain, and a witness inside $\mathbb{C}$ need not meet a
prescribed subdomain, so a statement about $\mathbb{C}$ would not restrict to
one about $U$.  And the Sobolev exponent is the endpoint $q=2/(2-\alpha)$ and
nothing smaller: $(2-\alpha)q=2$ is what makes the Jacobian factor cancel in
the scaling identity (59), and for smaller $q$ the assertion is false, which is
why the positive half (10) of Theorem 1 carries a strict inequality.

The properties of $f$ are taken as hypotheses, as they are in the other nodes
of the Theorem 1 chain; they are properties of the inverse of (54), and are
discharged where this theorem is applied.
\end{theorem}
\begin{proof}
Throughout, $e_0=1$ and $e_1=i$, and for $h\in\mathbb{C}$ we write
$h^{(0)}=\operatorname{Re}h$, $h^{(1)}=\operatorname{Im}h$, so that
$h=h^{(0)}e_0+h^{(1)}e_1$.  A word $\nu=(j_1,\dots,j_m)$ in $\{0,1\}$ has
length $|\nu|=m$ and $D^\nu=\partial_{j_1}\cdots\partial_{j_m}$ as in
\noderef{PartialWord}, so that $D^{(j,\nu)}=\partial_j D^\nu$.  Finally
$c_\nu(h)=h^{(j_1)}\cdots h^{(j_m)}$.

\emph{Step 0: elementary consequences of the hypotheses.}

\noindent(0a) $d=k-1+\alpha$.  By \noderef{KAlphaEqDExpAddOne} the right-hand
side of (9) equals $d+1$, and by hypothesis it equals $k+\alpha$; hence
$d=k+\alpha-1$.

\noindent(0b) $d>0$, since $k\ge1$ and $\alpha>0$ give $d\ge\alpha>0$.

\noindent(0c) $f(0)=0$.  Apply (H3) with $t=0$ and any $z$: $f(0)=0^{d}f(z)=0$,
because $d\neq0$ by (0b).

\noindent(0d) $f$ is not identically zero: $f(H(1))=1$ by (H2).

\noindent(0e) $\alpha<1$ if and only if $d<k$, and $\alpha=1$ if and only if
$d=k$; both by (0a).

\noindent(0f) $q\ge1$ and $(2-\alpha)q=2$.  Indeed $0<\alpha\le1$ gives
$1\le2-\alpha<2$, so $q=2/(2-\alpha)$ satisfies $q>1$ and
$(2-\alpha)q=2$.

\noindent(0g) If $\alpha<1$ then $d-j\neq0$ for every integer $0\le j\le k$.
By (0a), $k-1<d<k$; so for $j\le k-1$ we have $d-j\ge d-(k-1)=\alpha>0$, and
$d-k=\alpha-1<0$.

\emph{Step 1: three standard facts, recorded for use below.}

\noindent(1a) \emph{Decomposition of a differential.}  Let $F$ be
differentiable at $w$, $L=\mathrm{d}F(w)$, $a=F_z(w)$, $b=F_{\bar z}(w)$.  By
\noderef{DZ} and \noderef{DZbar}, $a=\tfrac12(L[1]-iL[i])$ and
$b=\tfrac12(L[1]+iL[i])$, so $L[1]=a+b$ and $L[i]=i(a-b)$, and for
$h=h^{(0)}+ih^{(1)}$,
\[
  L[h]=h^{(0)}(a+b)+h^{(1)}i(a-b)=a\,h+b\,\bar h .
\]

\noindent(1b) \emph{Derivatives along a ray.}  Let $\Omega$ be open, let $g$ be
$C^\infty$ on $\Omega$, let $z_0\in\mathbb{C}$ and let $I$ be an interval of
positive reals with $tz_0\in\Omega$ for $t\in I$.  Put $\psi(t)=g(tz_0)$.  Then
$\psi$ is $C^\infty$ on $I$ and
\begin{equation}\tag{1}
  \psi^{(m)}(t)=\sum_{|\nu|=m}c_\nu(z_0)\,D^\nu g(tz_0)\qquad(t\in I,\ m\ge0).
\end{equation}
For $m=0$ this is the definition.  If it holds for $m$, then each $D^\nu g$ is
$C^\infty$ on $\Omega$ and, for any $C^1$ function $G$ on $\Omega$,
\[
  \frac{\mathrm{d}}{\mathrm{d}t}G(tz_0)=\mathrm{d}G(tz_0)[z_0]
  =\sum_{j\in\{0,1\}}z_0^{(j)}\,(\partial_jG)(tz_0)
\]
by the chain rule and the $\mathbb{R}$-linearity of $\mathrm{d}G(tz_0)$ in the
direction.  Applying this to each $G=D^\nu g$ and using
$\partial_jD^\nu=D^{(j,\nu)}$ gives (1) for $m+1$.

\noindent(1c) \emph{Partial derivatives and iterated Fréchet derivatives.}  Let
$\Omega$ be open and $g$ of class $C^m$ on $\Omega$.  Since $\Omega$ is open,
the iterated Fréchet derivative of $g$ within $\Omega$ coincides on $\Omega$
with the iterated Fréchet derivative of $g$, which we write
$\mathrm{d}^mg(x)$; it is the symmetric $m$-linear map characterized by the
recursion $\mathrm{d}^{m+1}g(x)[v_0,v_1,\dots,v_m]
=\bigl(\mathrm{d}(\mathrm{d}^mg)(x)[v_0]\bigr)[v_1,\dots,v_m]$, with
$\mathrm{d}^0g(x)=g(x)$.  Comparing this recursion with
$D^{(j,\nu)}=\partial_jD^\nu$ gives, for every word $\nu=(j_1,\dots,j_m)$,
\begin{equation}\tag{2}
  D^\nu g(x)=\mathrm{d}^mg(x)[e_{j_1},\dots,e_{j_m}]\qquad(x\in\Omega),
\end{equation}
and consequently, $\|e_j\|=1$,
\begin{equation}\tag{3}
  \bigl|D^\nu g(x)-D^\nu g(y)\bigr|
  \le\bigl\|\mathrm{d}^mg(x)-\mathrm{d}^mg(y)\bigr\|\qquad(x,y\in\Omega).
\end{equation}
Also, $g$ of class $C^m$ on $\Omega$ makes each $D^\nu g$ with $|\nu|=m$
continuous on $\Omega$, and each $D^\nu g$ with $|\nu|=m-1$ of class $C^1$ on
$\Omega$.

\emph{Step 2: the radial vanishing lemma.}  Let $m\ge0$, let $R>0$, let
$|z_0|=1$, and suppose
\[
  D^\nu f(tz_0)=0\qquad\text{for every }|\nu|=m\text{ and every }0<t<R .
\]
Suppose moreover that $d-j\neq0$ for every integer $0\le j\le m-1$.  Then
$f(z_0)=0$.

Indeed, put $\psi(t)=f(tz_0)$ for $t\in(0,R)$.  By (H1) and (1b) with
$\Omega=\mathbb{C}-\{0\}$, which contains $tz_0$ for $0<t<R$, the function
$\psi$ is $C^\infty$ on $(0,R)$ and, by (1) and the hypothesis,
\[
  \psi^{(m)}(t)=\sum_{|\nu|=m}c_\nu(z_0)D^\nu f(tz_0)=0\qquad(0<t<R).
\]
On the other hand (H3) gives $\psi(t)=t^{d}f(z_0)$ for $t>0$, whence
$\psi^{(m)}(t)=d(d-1)\cdots(d-m+1)\,t^{\,d-m}f(z_0)$, the product being empty
and equal to $1$ when $m=0$.  Each factor is non-zero by assumption and
$t^{\,d-m}>0$, so $f(z_0)=0$.

\emph{Step 3: $f$ cannot vanish on the whole unit circle.}  Suppose
$f(z_0)=0$ for every $|z_0|=1$.  For $z\neq0$, (H3) with $t=|z|>0$ and the
point $z/|z|$ gives $f(z)=|z|^{d}f(z/|z|)=0$, and $f(0)=0$ by (0c).  So
$f\equiv0$, contradicting (0d).  Hence there is $z_0$ with $|z_0|=1$ and
$f(z_0)\neq0$.

\emph{Step 4: the $\alpha=1$ endgame.}  Assume $\alpha=1$, let $R>0$ with
$\overline{B}(0,R)\subseteq U$, and assume
\begin{equation}\tag{4}
  D^\nu f(z)=0\qquad\text{for every }|\nu|=k+1\text{ and every }
  0<|z|<R .
\end{equation}
We derive a contradiction.

\noindent(4a) The set $B=\{z:0<|z|<R\}$ is open, non-empty and connected: it is
the image of the connected set $(0,R)\times\mathbb{R}$ under the continuous map
$(r,\theta)\mapsto re^{i\theta}$.

\noindent(4b) By (H1), $f$ is $C^\infty$ on $B\subseteq\mathbb{C}-\{0\}$, and by
(4) every partial derivative of order $k+1$ of $f$ vanishes on $B$.  Hence
\noderef{VanishingDerivsIsPolynomial} applies on $B$ with $m=k$ and yields
coefficients $a_{ij}$, vanishing unless $i+j\le k$, with
\begin{equation}\tag{5}
  f(z)=\sum_{i+j\le k}a_{ij}z^i\bar z^{\,j}\qquad(z\in B).
\end{equation}

\noindent(4c) \emph{$f$ agrees on $B$ with the top-degree part.}  Put
$Q(z)=\sum_{i+j=k}a_{ij}z^i\bar z^{\,j}$.  By (0e), $d=k$.  Let $z\in B$ and
$0<s<1$; then $sz\in B$, and (H3) gives $f(sz)=s^{k}f(z)$.  Substituting (5) on
both sides,
\[
  \sum_{n=0}^{k}s^{n}\sum_{i+j=n}a_{ij}z^i\bar z^{\,j}
  =s^{k}\sum_{i+j\le k}a_{ij}z^i\bar z^{\,j}
  \qquad(0<s<1).
\]
Both sides are polynomials in $s$ of degree at most $k$ which agree at the
infinitely many points of $(0,1)$, so they agree coefficientwise.  Comparing
the coefficients of $s^{k}$ and using (5),
\begin{equation}\tag{6}
  Q(z)=\sum_{i+j=k}a_{ij}z^i\bar z^{\,j}
  =\sum_{i+j\le k}a_{ij}z^i\bar z^{\,j}=f(z)\qquad(z\in B).
\end{equation}

\noindent(4d) \emph{$f=Q$ on $\mathbb{C}-\{0\}$.}  Note $Q(tz)=t^{k}Q(z)$ for
$t>0$, each monomial having $i+j=k$ and $t$ being real.  Let $z\neq0$ and put
$t=R/(2|z|)>0$, so that $0<|tz|=R/2<R$ and $tz\in B$.  By (H3),
$f(tz)=t^{k}f(z)$, so by (6)
\[
  f(z)=t^{-k}f(tz)=t^{-k}Q(tz)=t^{-k}t^{k}Q(z)=Q(z).
\]

\noindent(4e) \emph{The identity on the unit circle.}  If $\xi\neq0$ then
$H(\xi)\neq0$: otherwise (H2) and (0c) give $\xi=f(H(\xi))=f(0)=0$.  Let
$|\xi|=1$.  Then $H(\xi)\neq0$, so by (H2) and (4d),
\begin{equation}\tag{7}
  \xi=f(H(\xi))=Q(H(\xi))
  =\sum_{m=0}^{k}a_{m,k-m}H(\xi)^{m}\overline{H(\xi)}^{\,k-m}.
\end{equation}
Write $\xi=e^{i\theta}$, which is possible for every $\xi$ of modulus $1$.  The
first clause of \noderef{ExtremalPolarForm} at $r=1$ gives
\[
  H(\xi)=1^{1/d}\bigl(e^{i\theta}+\varepsilon e^{-3i\theta}\bigr)
  =\xi+\varepsilon\xi^{-3},
\]
where $\varepsilon=\epsilon_2(1,p)$ is the real coefficient of
\noderef{EpsCoef}, and $e^{-3i\theta}=\xi^{-3}$.  Since $|\xi|=1$ we have
$\bar\xi=\xi^{-1}$, so $\overline{\xi^{-3}}=\xi^{3}$ and, $\varepsilon$ being
real,
\[
  \overline{H(\xi)}=\xi^{-1}+\varepsilon\xi^{3}.
\]
Substituting into (7),
\begin{equation}\tag{8}
  \xi=\sum_{m=0}^{k}a_{m,k-m}
    \bigl(\xi+\varepsilon\xi^{-3}\bigr)^{m}
    \bigl(\xi^{-1}+\varepsilon\xi^{3}\bigr)^{k-m}
  \qquad(|\xi|=1).
\end{equation}

\noindent(4f) \emph{Continuation of (8) to all $\xi\neq0$.}  For $\xi\neq0$,
\[
  \bigl(\xi+\varepsilon\xi^{-3}\bigr)^{m}=\xi^{-3m}\bigl(\xi^{4}+\varepsilon
  \bigr)^{m},\qquad
  \bigl(\xi^{-1}+\varepsilon\xi^{3}\bigr)^{k-m}
  =\xi^{-(k-m)}\bigl(1+\varepsilon\xi^{4}\bigr)^{k-m},
\]
so multiplying the $m$-th summand of (8) by $\xi^{3k}$ produces
$a_{m,k-m}\xi^{\,3k-3m-(k-m)}(\xi^{4}+\varepsilon)^{m}
(1+\varepsilon\xi^{4})^{k-m}$ with
$3k-3m-(k-m)=2(k-m)\ge0$.  Hence, setting
\[
  P_1(X)=X^{3k+1},\qquad
  P_2(X)=\sum_{m=0}^{k}a_{m,k-m}X^{2(k-m)}
    \bigl(X^{4}+\varepsilon\bigr)^{m}
    \bigl(1+\varepsilon X^{4}\bigr)^{k-m},
\]
both $P_1$ and $P_2$ are polynomials, and multiplying (8) by $\xi^{3k}$ gives
$P_1(\xi)=P_2(\xi)$ for every $\xi$ with $|\xi|=1$.  The unit circle is
infinite, $\theta\mapsto e^{i\theta}$ being injective on $[0,2\pi)$, so the
polynomial $P_1-P_2$ has infinitely many roots and is therefore the zero
polynomial.  Hence $P_1(\xi)=P_2(\xi)$ for every $\xi\in\mathbb{C}$, and
dividing by $\xi^{3k}\neq0$ and undoing the two displays above,
\begin{equation}\tag{9}
  \xi=\sum_{m=0}^{k}a_{m,k-m}
    \bigl(\xi+\varepsilon\xi^{-3}\bigr)^{m}
    \bigl(\xi^{-1}+\varepsilon\xi^{3}\bigr)^{k-m}
  \qquad(\xi\neq0).
\end{equation}

\noindent(4g) \emph{The contradiction.}  By \noderef{ExtremalEpsBound},
$\varepsilon=0$ if and only if $p=2$; since $p\neq2$ we have
$\varepsilon\neq0$.  Also $k\ge1$.  So \noderef{HomogPolyModFour} applies with
the coefficient $\varepsilon$, the degree $k$ and the coefficient family
$a_{ij}$, whose hypothesis is exactly (9), and it yields a contradiction.  This
completes Step 4.

\emph{Step 5: $f\notin C^{k-1+\alpha}_{\mathrm{loc}}(U)$.}  Suppose to the
contrary that $f\in C^{k-1+\alpha}_{\mathrm{loc}}(U)$ in the sense of
\noderef{MemHolderSmallOn} at the index $k-1$.  Since $0\in U$ and $U$ is open,
fix $\rho>0$ with $K=\overline{B}(0,\rho)\subseteq U$; $K$ is compact.  By
\noderef{MemHolderSmallOn} one of two cases holds.

\noindent\emph{Case A: $\alpha<1$.}  Then $f$ is of class $C^{k-1}$ on $U$ and,
for every compact $K'\subseteq U$ and every $\delta>0$, there is $\rho'>0$ with
\begin{equation}\tag{10}
  \bigl\|\mathrm{d}^{k-1}f(x)-\mathrm{d}^{k-1}f(y)\bigr\|
  \le\delta|x-y|^{\alpha}\qquad(x,y\in K',\ |x-y|<\rho'),
\end{equation}
where we have used (1c) to replace the derivative within $U$ by the derivative
at the point, $U$ being open.

Fix a word $\nu$ with $|\nu|=k-1$.

First, $D^\nu f(0)=0$.  By (1c), $D^\nu f$ is continuous on $U$.  Pick any $z$
with $0<|z|\le\rho$; then $tz\in K\subseteq U$ for $0<t\le1$, and (H4) gives
\[
  D^\nu f(tz)=t^{\,d-(k-1)}D^\nu f(z)=t^{\alpha}D^\nu f(z)\longrightarrow0
  \qquad(t\to0^+),
\]
using (0a) and $\alpha>0$.  Since $tz\to0$ and $D^\nu f$ is continuous at
$0\in U$, we get $D^\nu f(0)=0$.

Next, let $\delta>0$ and let $\rho'>0$ be as in (10) for $K'=K$.  Fix $z$ with
$0<|z|\le\rho$ and let $0<t\le1$ satisfy $t|z|<\rho'$.  Then $x=tz$ and $y=0$
lie in $K$ and $|x-y|=t|z|<\rho'$, so by (10) and (3),
\[
  \bigl|D^\nu f(tz)\bigr|=\bigl|D^\nu f(tz)-D^\nu f(0)\bigr|
  \le\delta\,(t|z|)^{\alpha}.
\]
Combining with $D^\nu f(tz)=t^{\alpha}D^\nu f(z)$ and dividing by
$t^{\alpha}>0$,
\[
  \bigl|D^\nu f(z)\bigr|\le\delta|z|^{\alpha}.
\]
As $\delta>0$ was arbitrary, $D^\nu f(z)=0$.  Thus
\begin{equation}\tag{11}
  D^\nu f(z)=0\qquad\text{for every }|\nu|=k-1\text{ and }0<|z|\le\rho .
\end{equation}

Now let $|z_0|=1$.  For $0<t<\rho$ we have $0<|tz_0|=t\le\rho$, so (11) gives
$D^\nu f(tz_0)=0$ for every $|\nu|=k-1$.  By (0g) the numbers $d-j$ are
non-zero for $0\le j\le k-2$, so Step 2 with $m=k-1$ and $R=\rho$ gives
$f(z_0)=0$.  This holds for every $z_0$ on the unit circle, contradicting
Step 3.

\noindent\emph{Case B: $\alpha=1$.}  Then, by \noderef{MemHolderSmallOn} at the
index $k-1$ and $\alpha=1$, the function $f$ is of class $C^{(k-1)+1}=C^{k}$ on
$U$, using $k\ge1$.

Fix a word $\nu$ with $|\nu|=k-1$.  By (1c), $D^\nu f$ is of class $C^1$ on
$U$; in particular it is continuous on $U$ and differentiable at $0$.  Exactly
as in Case A --- using (H4) with $d-(k-1)=\alpha=1$ --- we get $D^\nu f(0)=0$.
Put $a=\partial_z(D^\nu f)(0)$ and $b=\partial_{\bar z}(D^\nu f)(0)$.  By
differentiability at $0$, $D^\nu f(0)=0$ and (1a),
\begin{equation}\tag{12}
  D^\nu f(\zeta)=a\zeta+b\bar\zeta+o(|\zeta|)\qquad(\zeta\to0).
\end{equation}
Fix $z\neq0$.  For all small $t>0$ we have $tz\in U$, and (H4) with
$|\nu|=k-1$ and $d-(k-1)=1$ gives $D^\nu f(tz)=t\,D^\nu f(z)$.  Hence by (12)
with $\zeta=tz$,
\[
  t\,D^\nu f(z)=a(tz)+b\overline{(tz)}+o(t)=t\bigl(az+b\bar z\bigr)+o(t) ,
\]
and dividing by $t$ and letting $t\to0^+$,
\begin{equation}\tag{13}
  D^\nu f(z)=a_\nu z+b_\nu\bar z\qquad(z\neq0),
\end{equation}
where we now write $a_\nu,b_\nu$ for the constants $a,b$ attached to $\nu$.

Let $\nu'$ be any word with $|\nu'|=k+1$.  Since $k\ge1$, $|\nu'|\ge2$, so we
may write $\nu'=(j_1,j_2,\nu)$ with $|\nu|=k-1$, and then
$D^{\nu'}f=\partial_{j_1}\partial_{j_2}(D^\nu f)$.  On the open set
$\mathbb{C}-\{0\}$ the function $D^\nu f$ agrees by (13) with the
$\mathbb{R}$-linear map $z\mapsto a_\nu z+b_\nu\bar z$, whose first partial
derivatives are the constants $a_\nu e_{j}+b_\nu\overline{e_{j}}$; hence
$\partial_{j_2}(D^\nu f)$ is constant on $\mathbb{C}-\{0\}$ and
$D^{\nu'}f=\partial_{j_1}\partial_{j_2}(D^\nu f)=0$ there.

Finally choose $R>0$ with $\overline{B}(0,R)\subseteq U$.  Then (4) holds, and
Step 4 gives a contradiction.

In both cases we have reached a contradiction, so
$f\notin C^{k-1+\alpha}_{\mathrm{loc}}(U)$.

\emph{Step 6: $f\notin W^{k+1,q}_{\mathrm{loc}}(U)$.}  Suppose to the contrary
that $f\in W^{k+1,q}_{\mathrm{loc}}(U)$ in the sense of \noderef{MemWLocOn}.
Let $(D_w)_w$ be the corresponding family, indexed by words, so that
$D_{()}=f$, each $D_w$ with $|w|\le k+1$ lies in $L^{q}_{\mathrm{loc}}(U)$ in
the sense of \noderef{MemLpLocOn}, and $D_{w\frown(j)}$ is the distributional
$j$-th partial derivative of $D_w$ on $U$ in the sense of
\noderef{IsWeakPartialOn} whenever $|w|<k+1$, where $w\frown(j)$ denotes $w$
with the letter $j$ appended on the right.  Put
$\Omega'=U-\{0\}$, an open set on which $f$ is $C^\infty$ by (H1).

\noindent(6a) \emph{Identification of the weak derivatives.}  We claim that for
every word $w$ with $|w|\le k+1$,
\begin{equation}\tag{14}
  D_w=D^{\overleftarrow{w}}f\quad\text{almost everywhere on }\Omega',
\end{equation}
where $\overleftarrow{w}$ is $w$ reversed.  We induct on $|w|$.  For $|w|=0$
both sides are $f$.  Assume (14) for all words of length $s<k+1$ and let
$|w|=s$, $j\in\{0,1\}$; we prove it for $w\frown(j)$, whose reverse is
$(j,\overleftarrow{w})$.

Put $u=D^{\overleftarrow{w}}f$, which is $C^\infty$ on $\Omega'$.  Every test
function for $\Omega'$ is a test function for $U$, its support being contained
in $\Omega'\subseteq U$; so, restricting the integration-by-parts identity of
\noderef{IsWeakPartialOn} to such test functions, $D_{w\frown(j)}$ is the
distributional $j$-th partial of $D_w$ on $\Omega'$.  Both integrals in that
identity are unchanged if $D_w$ is replaced by a function agreeing with it
almost everywhere on $\Omega'$, the integrands being supported in $\Omega'$; so
by the inductive hypothesis $D_{w\frown(j)}$ is also the distributional $j$-th
partial of $u$ on $\Omega'$, for each of the two letters $j$.

We may now apply \noderef{WeakPartialEqFDerivAE} on the open set $\Omega'$ to
the function $u$ and to the pair of candidate derivatives
$j\mapsto D_{w\frown(j)}$.  Its hypotheses hold: $u$ is continuous on
$\Omega'$, hence bounded on compact subsets of $\Omega'$ and therefore in
$L^1_{\mathrm{loc}}(\Omega')$; each $D_{w\frown(j)}$ lies in
$L^{q}_{\mathrm{loc}}(U)$ with $|w\frown(j)|=s+1\le k+1$, hence in
$L^1_{\mathrm{loc}}(\Omega')$, because compact subsets of $\Omega'$ are
compact subsets of $U$ of finite measure and $q\ge1$ by (0f), so Hölder's
inequality bounds the $L^1$ norm on such a set by the $L^{q}$ norm; the two
integration-by-parts identities were just established; and $u$, being
$C^\infty$ on $\Omega'$, is differentiable at every point of $\Omega'$.  The
conclusion is that, almost everywhere on $\Omega'$,
\[
  D_{w\frown(j)}=\partial_j u=\partial_j D^{\overleftarrow{w}}f
  =D^{(j,\overleftarrow{w})}f=D^{\overleftarrow{w\frown(j)}}f ,
\]
which is (14) for $w\frown(j)$.  This proves the claim.

\noindent(6b) \emph{Finiteness of the integral.}  Fix $R>0$ with
$\overline{B}(0,R)\subseteq U$ and fix a word $\nu$ with $|\nu|=k+1$; put
$F=D^\nu f$.  By (6a) applied to $w=\overleftarrow{\nu}$, we have $F=D_{w}$
almost everywhere on $\Omega'$, hence almost everywhere on $B(0,R)$, the origin
being a null set.  Since $\overline{B}(0,R)$ is a compact subset of $U$ and
$D_{w}\in L^{q}_{\mathrm{loc}}(U)$,
\begin{equation}\tag{15}
  \int_{B(0,R)}|F|^{q}\,\mathrm{d}\sigma
  =\int_{B(0,R)}|D_{w}|^{q}\,\mathrm{d}\sigma<\infty .
\end{equation}
Note also that $F$ is continuous on the open set $B=\{0<|z|<R\}$ by (H1), hence
measurable.

\noindent(6c) \emph{The scaling identity (59) and its integrated form.}  By
(H4) with $|\nu|=k+1$ and (0a),
$d-(k+1)=\alpha-2$, so
\begin{equation}\tag{16}
  F(z)=t^{\,2-\alpha}F(tz)\qquad(z\neq0,\ t>0).
\end{equation}
Let $0<t\le1$.  Using (16), then $(2-\alpha)q=2$ from (0f), then the change of
variables $\zeta=tz$, whose Jacobian is $t^{2}$,
\[
  \int_{B(0,R)}|F(z)|^{q}\,\mathrm{d}\sigma(z)
  =t^{(2-\alpha)q}\int_{B(0,R)}|F(tz)|^{q}\,\mathrm{d}\sigma(z)
  =t^{2}\,t^{-2}\int_{B(0,tR)}|F(\zeta)|^{q}\,\mathrm{d}\sigma(\zeta)
  =\int_{B(0,tR)}|F|^{q}\,\mathrm{d}\sigma ,
\]
the first equality holding because (16) is valid at every $z\neq0$, that is
almost everywhere.  This is the paper's identity
$\iint_{|z|<R}|D^\nu f|^{q}=\iint_{|z|<tR}|D^\nu f|^{q}$, and it is exactly
here that the endpoint exponent is used: for $q<2/(2-\alpha)$ the two powers of
$t$ do not cancel.

\noindent(6d) \emph{$F$ vanishes.}  By (15) the function $|F|^{q}$ is integrable
on $B(0,R)$, and the sets $B(0,tR)$ decrease to $\{0\}$ as $t\to0^+$, a null
set; so by dominated convergence
$\int_{B(0,tR)}|F|^{q}\to0$.  By (6c) the left-hand side of the display is
independent of $t$, so
\[
  \int_{B(0,R)}|F|^{q}\,\mathrm{d}\sigma=0 ,
\]
whence $F=0$ almost everywhere on $B(0,R)$.  Since $F$ is continuous on the
open set $B=\{0<|z|<R\}$ and the complement of a null set is dense, $F=0$ at
every point of $B$.  As $\nu$ with $|\nu|=k+1$ was arbitrary,
\begin{equation}\tag{17}
  D^\nu f(z)=0\qquad\text{for every }|\nu|=k+1\text{ and every }0<|z|<R .
\end{equation}

\noindent(6e) \emph{The contradiction, in two branches.}  If $\alpha<1$, let
$|z_0|=1$; for $0<t<R$ we have $0<|tz_0|=t<R$, so (17) gives
$D^\nu f(tz_0)=0$ for every $|\nu|=k+1$, and by (0g) the numbers $d-j$ are
non-zero for $0\le j\le k$, so Step 2 with $m=k+1$ and this $R$ gives
$f(z_0)=0$.  This holds for every $z_0$ on the unit circle, contradicting
Step 3.  If $\alpha=1$, then (17) is exactly (4) and Step 4 gives the
contradiction directly.

Hence $f\notin W^{k+1,q}_{\mathrm{loc}}(U)$, and with Step 5 the theorem is
proved.
\end{proof}
